\section{Discussion}
\label{sec:discussion}

We have presented a comprehensive zero-shot variant effect atlas for ten
spaceflight radiation-response genes, scoring 215,001 variants spanning coding,
non-coding, and indel categories using the Evo2 genomic foundation model. Our
three-layer validation framework demonstrates that Evo2 zero-shot predictions
are competitive with supervised tools for clinical variant classification and
offer unique capabilities for non-coding and indel variant interpretation.

\subsection*{Foundation models as complementary VEP tools}

Evo2's zero-shot approach occupies a distinct niche in the variant effect
prediction landscape. Unlike AlphaMissense, which achieves high accuracy for
missense variants through AlphaFold2 structural priors, or CADD, which
integrates dozens of annotation tracks including ClinVar itself, Evo2 derives
predictions solely from sequence context without exposure to phenotypic labels.
This independence from training data confers two advantages: absence of
ClinVar circularity and the ability to score any variant type---including
promoter, intronic, and indel variants---through a single unified framework.

Our matched-variant analysis reveals that on the shared variant space
(predominantly missense), supervised tools---particularly AlphaMissense---retain
an accuracy advantage. This is expected: missense prediction is a well-defined
protein-level problem for which structure-informed models are specifically
optimised. However, the orthogonality analysis demonstrates that Evo2 captures
complementary information (Spearman $\rho = 0.41$--$0.83$ vs.\ CADD;
$\rho = 0.05$--$0.45$ vs.\ SpliceAI), and the ensemble of Evo2 + CADD
consistently matches or exceeds either tool alone. The TP53 ensemble improvement
($+0.38\%$ AUROC) is particularly notable given the already high baseline
performance, suggesting that the two tools identify partially non-overlapping
sets of true positives.

\subsection*{Non-coding and indel variants: an unmet need}

The most compelling evidence for foundation-model VEP comes from non-coding
variants. For the TERT promoter, Evo2 was the only tool with a significant
positive correlation with MPRA functional data ($\rho = +0.267$,
$p = 3.8 \times 10^{-14}$), while CADD, SpliceAI, and ncER showed
non-significant or negatively correlated predictions. This result highlights a
genuine gap in the current VEP toolkit: existing non-coding scores either lack
variant-level resolution (ncER provides regional scores) or are not designed for
promoter-regulatory logic (SpliceAI targets splice junctions). Evo2's ability to
capture promoter-level sequence grammar, trained implicitly through genome-wide
likelihood optimisation, fills this gap.

Similarly, the high frameshift indel AUROCs (mean 0.977 across six genes) demonstrate
that the log-likelihood framework naturally extends to length-altering variants.
AlphaMissense and REVEL cannot score indels by design, and CADD requires
separate pre-computed indel annotations. Evo2's unified scoring of all variant
types from a single model forward pass represents a practical advantage for
clinical pipelines that must process mixed variant sets.

\subsection*{Mutation signatures: transversion bias, not radiation specificity}

Our initial observation that radiation-type mutations (C$>$A, G$>$T) receive
more damaging Evo2 scores than other mutation classes required careful control.
The transversion control analysis revealed that 8 of 10 genes showed no
significant difference between radiation-oxidative and other transversions after
Bonferroni correction, with only TP53 and DNMT3A showing borderline effects
(rank-biserial $r_{bc} < 0.05$). The dominant signal is transversion versus
transition: purine$\leftrightarrow$pyrimidine substitutions cause greater
disruption to sequence grammar than purine$\leftrightarrow$purine or
pyrimidine$\leftrightarrow$pyrimidine changes, and Evo2 faithfully captures
this physical asymmetry.

This honest null result has practical implications: Evo2 is a sequence-disruption
detector, not a radiation exposure biomarker. However, the consistently greater
predicted impact of transversions is itself relevant to spaceflight risk
assessment, as the GCR mutational spectrum is enriched for transversions relative
to the endogenous mutational background.

\subsection*{Limitations}

Several limitations warrant discussion. First, Evo2's zero-shot accuracy, while
competitive, does not uniformly surpass supervised tools on their home turf---particularly
AlphaMissense for missense and CADD for ClinVar-trained categories. The appropriate
use case for Evo2 is as a complementary tool, not a replacement. Second, our
validation is constrained to genes with existing DMS, ClinVar, or MPRA data; the
six novel spaceflight genes (ATM, NFE2L2, CLOCK, MSTN, RAD51, TERT) have
variable annotation depth, and the predictions for CLOCK and MSTN (which lack
ClinVar pathogenic variants) remain experimentally unvalidated. Third, the TERT
C228T and C250T promoter hotspots---the two most clinically important TERT
variants---receive low Evo2 scores because they are gain-of-function mutations
that create de novo ETS binding sites rather than disrupting existing sequence
grammar. This reflects a fundamental limitation of loss-of-function-oriented
$\Delta$LL scoring for gain-of-function variants. Fourth, inference requires
GPU access (48\,GB VRAM minimum), limiting deployment in resource-constrained
clinical settings. Fifth, the indel validation is limited to ClinVar frameshift
classification; DMS-level functional data for indels would strengthen the
calibration.

\subsection*{Future directions}

The variant effect atlas presented here serves as a foundation for several
downstream applications. First, Evo2 scores could be incorporated into
ACMG/AMP PP3/BP4 evidence frameworks for variant interpretation, particularly
for non-coding and indel variants where existing tools provide limited coverage.
Second, ensemble approaches combining foundation-model scores with supervised
annotations warrant systematic evaluation across larger gene panels and diverse
population cohorts. Third, pre-computed Evo2 atlases for mission-critical gene
panels could support astronaut genomic risk assessment during mission planning.
Finally, as genomic foundation models continue to scale, the zero-shot VEP
paradigm may increasingly complement or substitute for supervised approaches
that require continual retraining on expanding clinical databases.
